% chapter/10_conclusion.tex
% Session 10 — Conclusion
% Compiled with pdflatex + bibtex (nonstopmode)

\section{Conclusion}
\label{sec:conclusion}

This paper introduces DCA-Trie, a dynamic context-aware constrained decoding framework
for knowledge graph question answering. The framework modifies only the constraint
oracle: a TypeOracle with two symbolic gates (range and type) replaces GCR's static
trie, filtering semantically irrelevant paths during beam search.

The TypeOracle reduces SIR from 1.000 to 0.145, an 85.5\% improvement in constraint
quality, while maintaining a false negative rate of 3.3\% and adding less than 1\% latency
overhead. These results demonstrate that purely symbolic gates, operating in $O(1)$ per
candidate triple, can achieve substantial search-space reduction without neural inference
or learned thresholds.

The paper's central finding is that constraint tightness and answer accuracy are
independent variables. DCA-Trie v1 filters 85.5\% of irrelevant paths, but Hits@1
drops from 91.6\% to 86.4\%. Four mechanisms govern this relationship: gold-path
exclusion, beam competition, type-inference noise, and the precision-recall trade-off.
The result holds across beam widths and is statistically significant ($p < 0.001$).
This finding challenges the implicit assumption in GCR and DoG that tightening the
search space necessarily improves reasoning.

The Semantic Irrelevance Ratio (SIR) provides the first metric that measures constraint
quality independently of answer accuracy. SIR enables systematic comparison of oracles
without conflating structural filtering with downstream performance. A system can have
excellent SIR (tight constraints) and poor Hits@1 (wrong answer), or poor SIR (loose
constraints) and excellent Hits@1 (correct answer). Before SIR, these two dimensions
were entangled.

The contributions are threefold. First, the TypeOracle demonstrates that symbolic gates
can replace embedding-based oracles with $O(1)$ complexity, no threshold tuning, and
deterministic behaviour. Second, SIR provides a diagnostic metric for constraint quality
that complements Hits@1. Third, the non-monotone finding establishes that structural
constraint quality and reasoning accuracy should be optimised independently, not assumed
to be correlated.

Future work has three promising directions. Fine-tuning the LLM on DCA-Trie's filtered
trie may teach the model to operate with reduced diversity, closing the Hits@1 gap. A
neural type classifier would improve the type gate's accuracy and reduce FNR. Adaptive
oracle thresholds that vary by hop depth could preserve diversity where it matters while
filtering where it does not.

The broader implication is that constrained decoding systems should report both accuracy
and constraint quality metrics. Hits@1 alone is insufficient: a system that achieves
high Hits@1 with high SIR may be over-constrained and fragile. SIR reveals the operating
point and exposes the trade-off that accuracy metrics alone cannot capture.